\documentclass[10pt,letterpaper]{article}

\usepackage[margin=0.78in]{geometry}
\usepackage[T1]{fontenc}
\usepackage[utf8]{inputenc}
\usepackage{lmodern}
\usepackage{microtype}
\usepackage{xcolor}
\usepackage{tabularx}
\usepackage{booktabs}
\usepackage{array}
\usepackage{enumitem}
\usepackage{amsmath}
\usepackage{hyperref}
\usepackage{fancyhdr}
\usepackage{titlesec}
\usepackage{parskip}

\definecolor{ink}{HTML}{18212F}
\definecolor{navy}{HTML}{183B5B}
\definecolor{blue}{HTML}{2F6F9F}
\definecolor{green}{HTML}{28785E}
\definecolor{slate}{HTML}{5B6573}
\definecolor{pale}{HTML}{F4F7FA}
\definecolor{line}{HTML}{D8DEE6}

\hypersetup{
    colorlinks=true,
    linkcolor=navy,
    urlcolor=blue
}

\pagestyle{fancy}
\fancyhf{}
\renewcommand{\headrulewidth}{0pt}
\renewcommand{\footrulewidth}{0.3pt}
\fancyfoot[L]{\small\color{slate}Auto Insurance Claim Probability Modeling}
\fancyfoot[R]{\small\color{slate}\thepage}

\titleformat{\section}
  {\large\bfseries\color{navy}}
  {}
  {0pt}
  {}
  [\vspace{-0.2em}{\color{line}\titlerule}]

\titleformat{\subsection}
  {\normalsize\bfseries\color{navy}}
  {}
  {0pt}
  {}

\setlist[itemize]{leftmargin=1.15em,itemsep=0.12em,topsep=0.2em}
\renewcommand{\arraystretch}{1.18}
\setlength{\tabcolsep}{5pt}
\newcolumntype{L}[1]{>{\raggedright\arraybackslash}p{#1}}
\newcolumntype{B}[1]{>{\raggedright\arraybackslash\bfseries}p{#1}}
\color{ink}

\newcommand{\summarybox}[3]{%
  \fcolorbox{line}{pale}{%
    \begin{minipage}[t][0.66in][t]{0.215\textwidth}
      \vspace{0.06in}
      {\scriptsize\color{slate}\MakeUppercase{#1}}\par
      {\Large\bfseries\color{navy}#2}\par
      {\scriptsize\color{slate}#3}
    \end{minipage}
  }%
}

\newcommand{\rocbar}[2]{%
  \begingroup
  \setlength{\fboxsep}{0pt}%
  \textcolor{#2}{\rule{#1\linewidth}{0.5em}}%
  \endgroup
}

\begin{document}

{\color{navy}\rule{\textwidth}{2.8pt}}
\vspace{0.35em}

{\LARGE\bfseries Auto Insurance Claim Probability Modeling}

\vspace{0.2em}
{\color{slate}A statistical regression workflow for policy-level auto insurance claim occurrence}

\vspace{0.9em}
\noindent
\summarybox{Policies}{58,592}{policy records}\hfill
\summarybox{Claim Rate}{6.40\%}{portfolio base rate}\hfill
\summarybox{Best ROC-AUC}{0.606}{R logistic model}\hfill
\summarybox{Top-Decile Lift}{1.68x}{highest-risk decile}

\section{Abstract}

This report analyzes auto insurance claim occurrence using a reproducible statistical modeling workflow. The response variable is a binary indicator of whether a policy generated a claim. The analysis begins with simple linear regression, extends to a multiple linear probability model, and compares both linear specifications with logistic regression.

The project is intentionally structured around interpretable models. Linear regression provides a transparent baseline, while logistic regression addresses the probability constraints required for a binary outcome. Holdout results show moderate risk segmentation: the R logistic regression model achieves a ROC-AUC of 0.606 and a top-decile lift of 1.68x.

\section{Data}

\begin{tabularx}{\textwidth}{>{\bfseries}p{1.65in}X}
\toprule
Item & Value \\
\midrule
Observation grain & One row per auto insurance policy \\
Rows & 58,592 \\
Columns & 41, including the target variable \\
Claim count & 3,748 \\
Claim rate & 6.40\% \\
Target variable & \texttt{claim\_status}: 1 if a claim occurred, 0 otherwise \\
Source handling & Raw Kaggle CSV excluded from the public repository; generated outputs and code are included \\
\bottomrule
\end{tabularx}

\section{Modeling Framework}

The modeling sequence follows a baseline-to-GLM progression commonly used in insurance analytics: start with an interpretable benchmark, add rating-style features, then use a model designed for binary outcomes.

\begin{tabularx}{\textwidth}{B{1.85in}L{1.75in}X}
\toprule
Component & Method & Role in the analysis \\
\midrule
Baseline & Simple linear regression & Uses \texttt{subscription\_length} as a one-variable reference model. \\
Linear probability model & Multiple linear regression & Adds policy, customer, region, and vehicle predictors while preserving a linear structure. \\
Binary response model & Logistic regression & Uses a binomial link so fitted probabilities remain within the valid \([0,1]\) interval. \\
\bottomrule
\end{tabularx}

\vspace{0.35em}
For the logistic model, the claim probability \(p_i\) for policy \(i\) is modeled as:

\[
\log\left(\frac{p_i}{1-p_i}\right)
= \beta_0 + \beta_1x_{i1} + \cdots + \beta_kx_{ik}.
\]

\section{Validation Strategy}

The portfolio is imbalanced because only 6.40\% of policies have a claim. Accuracy can therefore be misleading: a model can appear accurate by predicting no claim for almost every policy. The validation framework emphasizes ranking quality, probability error, and concentration of observed claims in the highest predicted-risk group.

\begin{tabularx}{\textwidth}{>{\bfseries}p{1.45in}X}
\toprule
Metric & Interpretation \\
\midrule
ROC-AUC & Measures whether policies with claims tend to receive higher predicted risk scores. \\
Brier score & Measures squared probability forecast error. Lower values indicate better probability estimates. \\
Precision and recall & Summarize the tradeoff at the portfolio base-rate threshold. \\
Top-decile lift & Compares observed claim frequency in the highest predicted-risk decile with the overall claim rate. \\
\bottomrule
\end{tabularx}

\section{Holdout Results}

\begin{tabularx}{\textwidth}{Xrrrrr}
\toprule
Model & ROC-AUC & Precision & Recall & Brier & Top-Decile Lift \\
\midrule
Simple linear baseline & 0.592 & 0.082 & 0.601 & 0.0595 & 1.42x \\
Multiple linear probability & 0.606 & 0.082 & 0.639 & 0.0594 & 1.63x \\
Logistic regression & 0.606 & 0.085 & 0.587 & 0.0594 & 1.68x \\
\bottomrule
\end{tabularx}

\vspace{0.45em}
\begin{tabularx}{\textwidth}{B{2.15in}Xr}
\toprule
Model & Relative ROC-AUC & ROC-AUC \\
\midrule
Simple linear baseline & \rocbar{0.592}{blue} & 0.592 \\
Multiple linear probability model & \rocbar{0.606}{green} & 0.606 \\
Logistic regression & \rocbar{0.606}{navy} & 0.606 \\
\bottomrule
\end{tabularx}

\section{Calibration Review}

The logistic model shows a generally increasing relationship between predicted and observed claim rates by risk decile. The highest predicted-risk decile has an observed claim rate of 10.7\%, compared with the portfolio base rate of 6.4\%.

\begin{center}
\small
\begin{tabular}{lrrrrrrrrrr}
\toprule
Decile & 1 & 2 & 3 & 4 & 5 & 6 & 7 & 8 & 9 & 10 \\
\midrule
Predicted & 3.0\% & 3.9\% & 4.4\% & 5.0\% & 5.6\% & 6.3\% & 7.1\% & 8.1\% & 9.2\% & 11.2\% \\
Observed & 2.7\% & 4.3\% & 4.8\% & 4.7\% & 6.0\% & 6.4\% & 7.0\% & 8.8\% & 8.5\% & 10.7\% \\
\bottomrule
\end{tabular}
\end{center}

\section{Discussion}

The main finding is not high predictive accuracy; it is moderate but measurable risk segmentation. Moving from the simple linear baseline to the logistic model raises top-decile lift from 1.42x to 1.68x. In practical terms, the highest predicted-risk decile contains policies with substantially higher observed claim frequency than the average test portfolio.

The multiple linear probability model performs similarly on ROC-AUC, but it can produce invalid probabilities outside the \([0,1]\) range. Logistic regression is therefore the cleaner probability model for this binary response setting, even when the ranking performance improvement is modest.

\section{Limitations}

\begin{itemize}
  \item The dataset does not include exposure duration in a form that supports a full claim frequency rate model.
  \item Claim severity and loss amount are not available, so this project does not estimate pure premium.
  \item Prior claims, coverage limits, deductibles, and richer territorial variables would likely improve model segmentation.
  \item The analysis should be viewed as a statistical modeling demonstration, not as a production pricing model.
\end{itemize}

\section{Reproducibility}

The project includes both R and Python implementations. The R workflow uses \texttt{lm()} for linear models and \texttt{glm(..., family = binomial)} for logistic regression. Generated outputs are stored as CSV files under \texttt{outputs/}, and public figures are stored under \texttt{figures/}.

\begin{tabularx}{\textwidth}{>{\bfseries}p{1.7in}X}
\toprule
Artifact & Location \\
\midrule
R model script & \texttt{R/claim\_probability\_models.R} \\
Python model script & \texttt{src/claim\_probability\_models.py} \\
Model metrics & \texttt{outputs/model\_metrics.csv}, \texttt{outputs/r\_model\_metrics.csv} \\
Report source & \texttt{docs/Auto\_Insurance\_Claim\_Probability\_Model\_Report.tex} \\
\bottomrule
\end{tabularx}

\end{document}
